\documentclass[11pt]{article}
\usepackage[portuguese]{babel}
\usepackage[utf8]{inputenc}
\usepackage[T1]{fontenc}
\usepackage{hyperref}
\hypersetup{colorlinks=true, linkcolor=blue, filecolor=magenta, urlcolor=cyan,}
\urlstyle{same}
\usepackage{amsmath}
\usepackage{amsfonts}
\usepackage{amssymb}
\usepackage[version=4]{mhchem}
\usepackage{stmaryrd}

\usepackage{tikz}
\usetikzlibrary{shapes, arrows}


\title{Para as Simulaçoes - TEX}

\author{}
\date{}


\begin{document}
\maketitle
\section*{Com base em: \href{https://stefvanbuuren.name/fimd/sec-evaluation.html}{https://stefvanbuuren.name/fimd/sec-evaluation.html}}
\begin{itemize}
  \item Medida que será imputada $Q$
  \item Vamos partir dos dados completos do Paraná (ou escolhe uma UF/Ano que tenha o menor percentual de missing
  \item Para os dados completos se calcula o \textit{true value} $Q$
  \item Podemos depois avaliar se iremos gerar uma população simulada
\end{itemize}

\begin{verbatim}
#install.packages("truncnorm")
library(truncnorm)

create.data <- function(beta = 1, sigma2 = 1, n = 50,
                        run = 1) {
  set.seed(seed = 12345)
  #x <- rnorm(n)
  x <- ceiling(rtruncnorm(n, a=12, b=45, mean=20, sd=10))
  #y <- beta * x + rnorm(n, sd = sqrt(sigma2))
   y <-  ceiling(rtruncnorm(n, a=12, b=85, mean=40, sd=25))
  data.frame(idade_mae = x, idade_pai = y)
}
dado.completo <- create.data(n=100)

hist(dado.completo$idade_mae)
hist(dado.completo$idade_pai)

boxplot(dado.completo$idade_mae,dado.completo$idade_pai)
\end{verbatim}

\subsection*{Métricas}
\textbf{Viés bruto (RB) e viés percentual (PB)}. O viés bruto da estimativa  $\overline{Q}$  é definido como a diferença entre o valor esperado da estimativa eo verdadeiro valor.

$$
RB = E(\overline{Q}) - Q
$$  O RB deve estar próximo de zero. O viés também pode ser expresso como viés percentual:

$$
PB = 100 \times \left| \frac{(E(\overline{Q}) - Q)}{Q} \right|
$$  Para um desempenho aceitável, utilizamos um limite superior para o PB de 5\%. (Demirtas, Freels e Yucel 2008)

\#\# Tradução para Português

\textbf{Erro quadrático médio (RMSE)}. O

$$
RMSE = \sqrt{(E(\overline{Q}) - Q)^2}
$$

é um compromisso entre viés e variância, e avalia  $\overline{Q}$   em termos de precisão e exatidão.

\subsection*{Processos sugeridos de simulação}
\textbf{Amostragem apenas}: Os passos básicos da simulação são: escolher\\
$Q$, coletar amostras  $Y(s)$, ajustar o modelo/imputação de dados completos, estimar  $\hat{Q}(s)$  e  $U(s)$  e calcular os resultados agregados sobre  $s$.

\textbf{Amostragem e mecanismo de dados ausentes combinados}: Os passos básicos da simulação são: escolher  $Q$, coletar amostras  $Y(s)$, gerar dados incompletos  $Y(s,t)_{obs}$, imputar, estimar  $\hat{Q}(s,t)$  e  $T(s,t)$  e calcular os resultados agregados sobre  $s$  e  $t$.

\textbf{Mecanismo de dados ausentes apenas} Os passos básicos da simulação são: escolher  $(\hat{Q}, U)$, gerar dados incompletos  $Y(t)_{obs}$, imputar, estimar  $(\overline{Q}, \overline{U})(t)$  e  $B(t)$  e calcular os resultados agregados sobre  $t$.

\section*{Do artigo \href{https://pmc.ncbi.nlm.nih.gov/articles/PMC8426774/}{https://pmc.ncbi.nlm.nih.gov/articles/PMC8426774/}}
\textbf{Viés Absoluto (|viés|) e Erro Quadrático Médio (RMSE)} de IM (imputação múltipla) medido sob as suposições de Dados Ausentes Completamente ao Acaso (MCAR), Dados Ausentes ao Acaso (MAR) e Não Ausentes ao Acaso (NMAR).

\subsection*{Viés Absoluto}
O viés absoluto é uma medida da precisão de um estimador, definido como a média das diferenças absolutas entre os valores estimados e os valores verdadeiros. A fórmula para o viés absoluto é:

$$
| B| = \frac{1}{n} \sum_{i=1}^{n} | \hat{Q}_i - Q_i |
$$

onde $\hat{Q}_i$ é o valor estimado e $Q_i$ é o valor verdadeiro e $n$ é a quantidade de amostras geradas (quantas vezes a amostragem foi repetida).

\subsection*{Erro Quadrático Médio (RMSE)}
O erro quadrático médio é uma medida que quantifica a diferença entre os valores previstos e os valores observados, calculado como a raiz quadrada da média dos quadrados das diferenças:

$$
RMSE = \sqrt{\frac{1}{n} \sum_{i=1}^{n} (\hat{Q}_i - Q_i)^2}
$$

\subsection*{Considerações sobre MCAR, MAR e NMAR}
\begin{itemize}
  \item \textbf{MCAR (Dados Ausentes Completamente ao Acaso)}: As observações ausentes não estão relacionadas a nenhuma variável observada ou não observada. Isso implica que o viés absoluto e o RMSE podem ser minimizados, pois a ausência de dados não introduz viés.
  \item \textbf{MAR (Dados Ausentes ao Acaso)}: As observações ausentes podem estar relacionadas a variáveis observadas, mas não às variáveis ausentes. O viés absoluto pode ser maior do que em MCAR, e o RMSE pode ser afetado dependendo da natureza das variáveis observadas.
  \item \textbf{NMAR (Não Ausentes ao Acaso)}: As observações ausentes estão relacionadas às próprias variáveis ausentes. Isso pode resultar em um viés absoluto significativo e um RMSE elevado, pois a ausência de dados pode introduzir um viés sistemático nas estimativas.
\end{itemize}

\subsection*{Adaptando a ideia}
\begin{itemize}
  \item Gerar um padrão univariado de dados ausentes
  \item Remover algumas observações para a variável Idade do Pai no conjunto de dados-completos, enquanto todas as outras variáveis são mantidas mantidas.
  \item Para MCAR: diferentes quantidades de observações para a variável (Idade do Pai)  deletadas aleatoriamente - considerar 0\%, 20\%, 40\%, 60\% e 80\% de dados ausentes.
  \item Para MAR : ordenar/classificar os dados completos de acordo com uma das variáveis completamente observadas (Idade da Mãe, variável $x$) e excluir os valores mais baixos (ou mais altos) dos casos da Idade do Pai em diferentes porcentagens dos valores ausentes para gerar o mecanismo MAR.
  \item Para o mecanismo NMAR: ordenar/classificar os dados completos pela própria variável ausente (Idade do Pai) e os valores da Idade do Pai são excluídos em cinco taxas diferentes
  \item Fazer isso  separadamente e independentemente para criar conjuntos de dados incompletos para os três mecanismos de ausência.
\end{itemize}

\textbf{Acho que pode valer a pena você pensar em gerar para diferentes mecanismos}

\section*{Com ajuda, tem-se os códigos para gerar os dados}
\begin{verbatim}
# Definindo a semente para reprodutibilidade
set.seed(123)

# Instalando e carregando o pacote necessário
# install.packages("truncnorm")
library(truncnorm)

create.data <- function(beta = 1, sigma2 = 1, n = 50, run = 1) {
  set.seed(seed = 123)
  x <- ceiling(rtruncnorm(n, a=12, b=45, mean=20, sd=10))
  y <- ceiling(rtruncnorm(n, a=12, b=85, mean=40, sd=25))
  data.frame(Idade_Mae = x, Idade_Pai = y)
}

dados <- create.data(n=100)

# Função para gerar dados MAR
gerar_mar <- function(dados, porcentagem) {
  # Ordenando os dados pela Idade da Mãe
  dados_ordenados <- dados[order(dados$Idade_Mae), ]
  
  # Calculando o número de observações a serem removidas
  n_remover <- floor(nrow(dados_ordenados) * porcentagem)
  
  # Removendo os n_remover menores valores da Idade do Pai
  dados_ordenados$Idade_Pai[1:n_remover] <- NA
  
  return(dados_ordenados)
}

# Função para gerar dados NMAR
gerar_nmar <- function(dados, porcentagem) {
  # Ordenando os dados pela Idade do Pai
  dados_ordenados <- dados[order(dados$Idade_Pai), ]
  
  # Calculando o número de observações a serem removidas
  n_remover <- floor(nrow(dados_ordenados) * porcentagem)
  
  # Removendo os n_remover maiores valores da Idade do Pai
  dados_ordenados$Idade_Pai[(nrow(dados_ordenados) - n_remover + 1):nrow(dados_ordenados)] <- NA
  
  return(dados_ordenados)
}

# Função para gerar dados MCAR
gerar_mcar <- function(dados, porcentagem) {
  # Copiando os dados para não alterar o original
  dados_mcar <- dados
  
  # Calculando o número de observações a serem removidas
  n_remover <- floor(nrow(dados_mcar) * porcentagem)
  
  # Selecionando aleatoriamente índices para remover
  indices_remover <- sample(1:nrow(dados_mcar), n_remover)
  
  # Removendo os valores da Idade do Pai nos índices selecionados
  dados_mcar$Idade_Pai[indices_remover] <- NA
  
  return(dados_mcar)
}

# Aplicando a função para diferentes porcentagens de dados ausentes para MAR
porcentagens_mar <- c(0.2, 0.4, 0.6, 0.8)  # Porcentagens de dados ausentes para MAR
resultados_mar <- setNames(lapply(porcentagens_mar, function(p) gerar_mar(dados, p)), 
                            paste0("MAR_", porcentagens_mar * 100))

# Aplicando a função para diferentes porcentagens de dados ausentes para NMAR
porcentagens_nmar <- c(0.2, 0.4, 0.6, 0.8)  # Porcentagens de dados ausentes para NMAR
resultados_nmar <- setNames(lapply(porcentagens_nmar, function(p) gerar_nmar(dados, p)), 
                              paste0("NMAR_", porcentagens_nmar * 100))

# Aplicando a função para diferentes porcentagens de dados ausentes para MCAR
porcentagens_mcar <- c(0.0, 0.2, 0.4, 0.6, 0.8)  # Porcentagens de dados ausentes para MCAR
resultados_mcar <- setNames(lapply(porcentagens_mcar, function(p) gerar_mcar(dados, p)), 
                                   paste0("MCAR_", porcentagens_mcar * 100))

View(resultados_mar$MAR_20)
View(resultados_mcar$MCAR_20)
View(resultados_nmar$NMAR_20)

\end{verbatim}

\section*{As simulações para MAR com p=0,2; 0,4; 0,6; 0,8}
Códigos feitos com ajuda de ferramentas (precisa ser verificado e modificado para os demais mecanismos!)

\begin{verbatim}
# Definindo a semente para reprodutibilidade
set.seed(123)

# Instalando e carregando os pacotes necessários
# install.packages("truncnorm")
# install.packages("dplyr")
# install.packages("knitr")
library(truncnorm)
library(dplyr)
library(knitr)

# Função para criar dados com Idade da Mãe e Idade do Pai -- Aqui vc pode colocar os dados do PARANÁ
create.data <- function(n = 50) {
  x <- ceiling(rtruncnorm(n, a=12, b=45, mean=20, sd=10))  # Idade da Mãe
  y <- ceiling(rtruncnorm(n, a=12, b=85, mean=40, sd=25))  # Idade do Pai
  data.frame(Idade_Mae = x, Idade_Pai = y)
}

# Função para gerar dados MAR com percentual de missing
gerar_mar <- function(dados, porcentagem) {
  dados_ordenados <- dados %>% arrange(Idade_Mae)
  n_remover <- floor(nrow(dados_ordenados) * porcentagem)
  dados_ordenados$Idade_Pai[1:n_remover] <- NA
  return(dados_ordenados)
}

# Função para imputar pela média
imputar_media <- function(dados) {
  media_idade_pai <- mean(dados$Idade_Pai, na.rm = TRUE)  # Calcula a média ignorando NAs
  dados$Idade_Pai[is.na(dados$Idade_Pai)] <- media_idade_pai  # Imputa a média nos NAs
  return(dados)
}

# Percentuais de missing a serem testados
percentuais_missing <- c(0.2, 0.4, 0.6, 0.8)
n_simulacoes <- 500

# Tabela para armazenar os resultados
resultados_tabela <- tibble(
  Percentual_Missing = numeric(),
  Media_Idade_Pai_Imputada = numeric(),
  Media_RB = numeric(),
  Media_RMSE = numeric()
)

# Simulação para cada percentual de missing
for (porcentagem in percentuais_missing) {
  
  medias_idades_imputadas <- numeric(n_simulacoes)  # Vetor para armazenar as médias
  todas_idades_pai <- numeric(n_simulacoes)          # Vetor para armazenar as idades do pai originais
  
  for (i in seq_len(n_simulacoes)) {
    dados <- create.data(n = 100)         # Cria os dados
    dados_mar <- gerar_mar(dados, porcentagem)          # Gera dados MAR com o percentual de missing atual
    dados_imputados <- imputar_media(dados_mar) # Imputa pela média
    
    medias_idades_imputadas[i] <- mean(dados_imputados$Idade_Pai) # Armazena a média das idades do pai imputadas
    todas_idades_pai[i] <- mean(dados$Idade_Pai)                 # Armazena a média das idades do pai originais
  }
  
  # Resultado final para o percentual atual: média das idades do pai imputadas em cada dado
  media_final <- mean(medias_idades_imputadas)
  
  # Resultados das métricas médias
  mean_rb <- mean(medias_idades_imputadas - todas_idades_pai)   # Viés Bruto médio
  mean_rmse <- sqrt(mean((medias_idades_imputadas - todas_idades_pai)^2)) # RMSE

  # Armazenando os resultados na tabela usando dplyr
  resultados_tabela <- resultados_tabela %>%
    add_row(
      Percentual_Missing = porcentagem * 100,
      Media_Idade_Pai_Imputada = media_final,
      Media_RB = mean_rb,
      Media_RMSE = mean_rmse
    )
}

# Exibindo a tabela final com os resultados usando kable
resultados_tabela %>%
  kable(caption = "Resultados da Simulação de Dados com Missingness")

\end{verbatim}

\section*{Para MCAR}
\begin{verbatim}
# Definindo a semente para reprodutibilidade
set.seed(123)

# Instalando e carregando os pacotes necessários
# install.packages("truncnorm")
# install.packages("dplyr")
# install.packages("knitr")
library(truncnorm)
library(dplyr)
library(knitr)

# Função para criar dados com Idade da Mãe e Idade do Pai
create.data <- function(n = 50) {
  x <- ceiling(rtruncnorm(n, a=12, b=45, mean=20, sd=10))  # Idade da Mãe
  y <- ceiling(rtruncnorm(n, a=12, b=85, mean=40, sd=25))  # Idade do Pai
  data.frame(Idade_Mae = x, Idade_Pai = y)
}

# Função para gerar dados MCAR com percentual de missing
gerar_mcar <- function(dados, porcentagem) {
  n_remover <- floor(nrow(dados) * porcentagem)
  indices_remover <- sample(1:nrow(dados), n_remover)  # Seleciona aleatoriamente os índices a serem removidos
  dados$Idade_Pai[indices_remover] <- NA  # Remove os valores da Idade do Pai nos índices selecionados
  return(dados)
}

# Função para imputar pela média
imputar_media <- function(dados) {
  media_idade_pai <- mean(dados$Idade_Pai, na.rm = TRUE)  # Calcula a média ignorando NAs
  dados$Idade_Pai[is.na(dados$Idade_Pai)] <- media_idade_pai  # Imputa a média nos NAs
  return(dados)
}

# Percentuais de missing a serem testados
percentuais_missing <- c(0, 0.2, 0.4, 0.6, 0.8)
n_simulacoes <- 500

# Tabela para armazenar os resultados
resultados_tabela <- tibble(
  Percentual_Missing = numeric(),
  Media_Idade_Pai_Imputada = numeric(),
  Media_RB = numeric(),
  Media_RMSE = numeric()
)

# Simulação para cada percentual de missing
for (porcentagem in percentuais_missing) {
  
  medias_idades_imputadas <- numeric(n_simulacoes)  # Vetor para armazenar as médias
  todas_idades_pai <- numeric(n_simulacoes)          # Vetor para armazenar as idades do pai originais
  
  for (i in seq_len(n_simulacoes)) {
    dados <- create.data(n = 100)         # Cria os dados
    dados_mcar <- gerar_mcar(dados, porcentagem)          # Gera dados MCAR com o percentual de missing atual
    dados_imputados <- imputar_media(dados_mcar) # Imputa pela média
    
    medias_idades_imputadas[i] <- mean(dados_imputados$Idade_Pai) # Armazena a média das idades do pai imputadas
    todas_idades_pai[i] <- mean(dados$Idade_Pai)                 # Armazena a média das idades do pai originais
  }
  
  # Resultado final para o percentual atual: média das idades do pai imputadas em cada dado
  media_final <- mean(medias_idades_imputadas)
  
  # Resultados das métricas médias
  mean_rb <- mean(medias_idades_imputadas - todas_idades_pai)   # Viés Bruto médio
  mean_rmse <- sqrt(mean((medias_idades_imputadas - todas_idades_pai)^2)) # RMSE

  # Armazenando os resultados na tabela usando dplyr
  resultados_tabela <- resultados_tabela %>%
    add_row(
      Percentual_Missing = porcentagem * 100,
      Media_Idade_Pai_Imputada = media_final,
      Media_RB = mean_rb,
      Media_RMSE = mean_rmse
    )
}

# Exibindo a tabela final com os resultados usando kable
resultados_tabela %>%
  kable(caption = "Resultados da Simulação de Dados com Missingness (MCAR)")

\end{verbatim}

\section*{Para NMAR}
\begin{verbatim}
# Definindo a semente para reprodutibilidade
set.seed(123)

# Instalando e carregando os pacotes necessários
# install.packages("truncnorm")
# install.packages("dplyr")
# install.packages("knitr")
library(truncnorm)
library(dplyr)
library(knitr)

# Função para criar dados com Idade da Mãe e Idade do Pai
create.data <- function(n = 50) {
  x <- ceiling(rtruncnorm(n, a=12, b=45, mean=20, sd=10))  # Idade da Mãe
  y <- ceiling(rtruncnorm(n, a=12, b=85, mean=40, sd=25))  # Idade do Pai
  data.frame(Idade_Mae = x, Idade_Pai = y)
}

# Função para gerar dados NMAR com percentual de missing
gerar_nmar <- function(dados, porcentagem) {
  dados_ordenados <- dados %>% arrange(Idade_Pai)  # Ordena pela Idade do Pai
  n_remover <- floor(nrow(dados_ordenados) * porcentagem)
  dados_ordenados$Idade_Pai[(nrow(dados_ordenados) - n_remover + 1):nrow(dados_ordenados)] <- NA  # Remove os maiores valores
  return(dados_ordenados)
}

# Função para imputar pela média
imputar_media <- function(dados) {
  media_idade_pai <- mean(dados$Idade_Pai, na.rm = TRUE)  # Calcula a média ignorando NAs
  dados$Idade_Pai[is.na(dados$Idade_Pai)] <- media_idade_pai  # Imputa a média nos NAs
  return(dados)
}

# Percentuais de missing a serem testados
percentuais_missing <- c(0.2, 0.4, 0.6, 0.8)
n_simulacoes <- 500

# Tabela para armazenar os resultados
resultados_tabela <- tibble(
  Percentual_Missing = numeric(),
  Media_Idade_Pai_Imputada = numeric(),
  Media_RB = numeric(),
  Media_RMSE = numeric()
)

# Simulação para cada percentual de missing
for (porcentagem in percentuais_missing) {
  
  medias_idades_imputadas <- numeric(n_simulacoes)  # Vetor para armazenar as médias
  todas_idades_pai <- numeric(n_simulacoes)          # Vetor para armazenar as idades do pai originais
  
  for (i in seq_len(n_simulacoes)) {
    dados <- create.data(n = 100)         # Cria os dados
    dados_nmar <- gerar_nmar(dados, porcentagem)          # Gera dados NMAR com o percentual de missing atual
    dados_imputados <- imputar_media(dados_nmar) # Imputa pela média
    
    medias_idades_imputadas[i] <- mean(dados_imputados$Idade_Pai) # Armazena a média das idades do pai imputadas
    todas_idades_pai[i] <- mean(dados$Idade_Pai)                 # Armazena a média das idades do pai originais
  }
  
  # Resultado final para o percentual atual: média das idades do pai imputadas em cada dado
  media_final <- mean(medias_idades_imputadas)
  
  # Resultados das métricas médias
  mean_rb <- mean(medias_idades_imputadas - todas_idades_pai)   # Viés Bruto médio
  mean_rmse <- sqrt(mean((medias_idades_imputadas - todas_idades_pai)^2)) # RMSE

  # Armazenando os resultados na tabela usando dplyr
  resultados_tabela <- resultados_tabela %>%
    add_row(
      Percentual_Missing = porcentagem * 100,
      Media_Idade_Pai_Imputada = media_final,
      Media_RB = mean_rb,
      Media_RMSE = mean_rmse
    )
}

# Exibindo a tabela final com os resultados usando kable
resultados_tabela %>%
  kable(caption = "Resultados da Simulação de Dados com Missingness (NMAR)")

\end{verbatim}

\noindent\rule{\textwidth}{0.5pt}
Claro! Um organograma pode ajudar a visualizar o processo de simulação para cada mecanismo de dados faltantes (MCAR, MAR e NMAR). Abaixo, descrevo como seria o organograma, e em seguida, forneço uma representação textual que você pode usar como referência para criar um gráfico visual.

\subsection*{Organograma do Processo de Simulação}
\begin{enumerate}
  \item \textbf{Início}

  \begin{itemize}
    \item Definir a semente para reprodutibilidade.
  \end{itemize}
  \item \textbf{Criar Dados}

  \begin{itemize}
    \item Gerar um conjunto de dados com duas variáveis: Idade da Mãe e Idade do Pai.
  \end{itemize}
  \item \textbf{Para Cada Percentual de Missing (20\%, 40\%, 60\%, 80\%)}

  \begin{itemize}
    \item \textbf{Simulação}:
    \begin{itemize}
      \item Repetir o processo para um número definido de simulações (ex: 500).
      \item \textbf{Gerar Dados Faltantes}:
      \begin{itemize}
        \item \textbf{MCAR}: Remover aleatoriamente uma porcentagem dos valores da Idade do Pai.
        \item \textbf{MAR}: Ordenar os dados pela Idade da Mãe e remover os menores valores da Idade do Pai.
        \item \textbf{NMAR}: Ordenar os dados pela Idade do Pai e remover os maiores valores.
      \end{itemize}
      \item \textbf{Imputação}:
      \begin{itemize}
        \item Imputar os valores faltantes pela média da Idade do Pai.
      \end{itemize}
      \item \textbf{Calcular Métricas}:
      \begin{itemize}
        \item Calcular a média das idades imputadas.
        \item Calcular o viés bruto (RB).
        \item Calcular o erro quadrático médio (RMSE).
      \end{itemize}
    \end{itemize}
  \end{itemize}
  \item \textbf{Armazenar Resultados}

  \begin{itemize}
    \item Armazenar as médias das idades imputadas, RB e RMSE em uma tabela.
  \end{itemize}
  \item \textbf{Exibir Resultados}

  \begin{itemize}
    \item Mostrar a tabela com os resultados formatados.
  \end{itemize}
  \item \textbf{Fim}

\end{enumerate}

\begin{verbatim}
[Início]
    |
    V
[Definir Semente]
    |
    V
[ Criar Dados ]
    |
    V
[Para Cada Percentual de Missing]
    |-----------------------|
    |                       |
    V                       V
[Simulação]              [Resultados]
    |                       |
    V                       |
[Repetir 500 vezes]       |
    |                       |
    V                       |
[Gerar Dados Faltantes]   |
    |                       |
    |-----------------------|
    |                       |
   /|\                     /|\
    |                       |
    V                       V
[MCAR]                  [Tabela com Resultados]
    |                       |
[Remover Aleatoriamente]   |
                           |
[MAR]                     [Exibir Resultados]
    |                       |
[Ordenar pela Idade da Mãe] 
    |
[Remover Menores Valores]
    |
[NMAR]
    |
[Ordenar pela Idade do Pai]
    |
[Remover Maiores Valores]
\end{verbatim}

\subsection*{Organograma em RMD}
\begin{verbatim}
# Carregar pacotes necessários
#install.packages("DiagrammeR")
library(DiagrammeR)
 
grViz("
digraph process {
  graph [layout = dot, rankdir = TB]
  
  node [shape = rectangle, style = filled, fillcolor = lightblue]
  
  start [label = 'Início', shape = ellipse]
  create_data [label = 'Dados Completos']
  for_each_percentage [label = 'Diferentes percentuais de dados faltantes']
  
  mcar [label = 'Gerar MCAR']
  mar [label = 'Gerar MAR']
  nmar [label = 'Gerar NMAR']
  
  imputar [label = 'Diferentes métodos de imputação']
  calcular_metricas [label = 'Calcular Métricas (RB, RMSE)']
  store_results [label = 'Armazenar Resultados']
  end [label = 'Fim', shape = ellipse]
  
  start -> create_data
  create_data -> for_each_percentage
  for_each_percentage -> mcar
  for_each_percentage -> mar
  for_each_percentage -> nmar
  
  mcar -> imputar
  mar -> imputar
  nmar -> imputar
  
  imputar -> calcular_metricas
  calcular_metricas -> store_results
  store_results -> end
}
")
\end{verbatim}

\section*{Organograma no latex}
\begin{verbatim}
\begin{tikzpicture}[
    % Definindo estilos para os nós
      node style/.style={rectangle, rounded corners, draw=black, minimum width=2cm, minimum height=1cm, align=center},
    startstop/.style={ellipse, draw=black},
    process/.style={rectangle, draw=black},
    arrow/.style={thick,->,>=stealth}
]

% Nós do organograma
\node[startstop] (start) {Início};
\node[process] (create) [below of=start] {Dados Completos};
\node[process] (percent) [below of=create] {Diferentes percentuais de \textit{missing}};

% Nós de decisão alinhados lado a lado
\node[process] (mcar) [below of=percent, xshift=-3cm] {Gerar MCAR};
\node[process] (mar) [below of=percent] {Gerar MAR};
\node[process] (nmar) [below of=percent, xshift=3cm] {Gerar NMAR};

\node[process] (impute) [below of=mcar, yshift=-1cm, xshift=3cm] {Diferentes métodos de Imputação};
\node[process] (metrics) [below of=impute] {Calcular Métricas (RB, RMSE)};
\node[process] (store) [below of=metrics] {Armazenar Resultados};
\node[startstop] (end) [below of=store] {Fim};

% Conectando os nós
\draw[arrow] (start) -- (create);
\draw[arrow] (create) -- (percent);
\draw[arrow] (percent) -- (mcar);
\draw[arrow] (percent) -- (mar);
\draw[arrow] (percent) -- (nmar);
\draw[arrow] (mar) -- (impute);
\draw[arrow] (mcar) -- (impute);
\draw[arrow] (nmar) -- (impute);
\draw[arrow] (impute) -- (metrics);
\draw[arrow] (metrics) -- (store);
\draw[arrow] (store) -- (end);

\end{tikzpicture}

\end{verbatim}

\noindent\rule{\textwidth}{0.5pt}
Referências

\href{https://stefvanbuuren.name/fimd/sec-evaluation.html}{https://stefvanbuuren.name/fimd/sec-evaluation.html}

\href{https://www.scielo.br/j/rbmet/a/RYYKg5VtRXXpSc9q3DRKPns/}{https://www.scielo.br/j/rbmet/a/RYYKg5VtRXXpSc9q3DRKPns/}

\href{https://www.youtube.com/watch?v=_pJX2SFoFtY}{https://www.youtube.com/watch?v=\_pJX2SFoFtY}

\href{https://analisemacro.com.br/econometria-e-machine-learning/mae-rmse-acc-f1-roc-r2-avaliacao-de-desempenho-de-modelos-preditivos/}{https://analisemacro.com.br/econometria-e-machine-learning/mae-rmse-acc-f1-roc-r2-avaliacao-de-desempenho-de-modelos-preditivos/}

\href{https://edisciplinas.usp.br/pluginfile.php/3043641/mod_resource/content/1/PROPRIEDADES}{https://edisciplinas.usp.br/pluginfile.php/3043641/mod\_resource/content/1/PROPRIEDADES} DOS ESTIMADORES.pdf

\href{https://pt.wikipedia.org/wiki/Erro_quadr%C3%A1tico_m%C3%A9dio}{https://pt.wikipedia.org/wiki/Erro\_quadrático\_médio}

\href{https://mariofilho.com/rmse-raiz-do-erro-quadratico-medio-em-machine-learning/}{https://mariofilho.com/rmse-raiz-do-erro-quadratico-medio-em-machine-learning/}

\href{https://pt.linkedin.com/advice/0/how-can-you-interpret-mean-squared-error-your-predictive-iifdf?lang=pt}{https://pt.linkedin.com/advice/0/how-can-you-interpret-mean-squared-error-your-predictive-iifdf?lang=pt}




\begin{tikzpicture}[
    % Definindo estilos para os nós
      node style/.style={rectangle, rounded corners, draw=black, minimum width=2cm, minimum height=1cm, align=center},
    startstop/.style={ellipse, draw=black},
    process/.style={rectangle, draw=black},
    arrow/.style={thick,->,>=stealth}
]

% Nós do organograma
\node[startstop] (start) {Início};
\node[process] (create) [below of=start] {Dados Completos};
\node[process] (percent) [below of=create] {Diferentes percentuais de \textit{missing}};

% Nós de decisão alinhados lado a lado
\node[process] (mcar) [below of=percent, xshift=-3cm] {Gerar MCAR};
\node[process] (mar) [below of=percent] {Gerar MAR};
\node[process] (nmar) [below of=percent, xshift=3cm] {Gerar NMAR};

\node[process] (impute) [below of=mcar, yshift=-1cm, xshift=3cm] {Diferentes métodos de Imputação};
\node[process] (metrics) [below of=impute] {Calcular Métricas (RB, RMSE)};
\node[process] (store) [below of=metrics] {Armazenar Resultados};
\node[startstop] (end) [below of=store] {Fim};

% Conectando os nós
\draw[arrow] (start) -- (create);
\draw[arrow] (create) -- (percent);
\draw[arrow] (percent) -- (mcar);
\draw[arrow] (percent) -- (mar);
\draw[arrow] (percent) -- (nmar);
\draw[arrow] (mar) -- (impute);
\draw[arrow] (mcar) -- (impute);
\draw[arrow] (nmar) -- (impute);
\draw[arrow] (impute) -- (metrics);
\draw[arrow] (metrics) -- (store);
\draw[arrow] (store) -- (end);

\end{tikzpicture}


\end{document}